\input{../preamble.tex}

\begin{document}

\section{Droite dirigée par un vecteur}
	
	\begin{multicols}{2}
	Considérons un point $A$ ainsi qu'un vecteur $v$ quelconques.
	Le translaté de $A$ par $v$ est noté $A+v$ et est représenté ci-contre.
	
		\begin{center}
		\begin{tikzpicture}[>=stealth, scale=1]
		\begin{axis}[xmin = -10, xmax=10, ymin=-10, ymax=10, axis x line=none, axis y line=none, axis line style=<->, xlabel={}, ylabel={}, ticks = none]
			\addplot[BLUE_E, mark=*, mark size = 1] (4,2) node[above] {$A$};	
			% v = (2,-1)
			\addplot[very thick, ->, PINK] (0,0) -- (axis cs:8,-4) node[midway, below] {$v$};
			% A+v
			\addplot[dashed, ->, GREY] (4,2) -- (axis cs:12,-2) node[midway, below] {$v$};
			\addplot[RED_E, mark=*, mark size = 1] (12,-2) node[above right] {$A+v$};	
		\end{axis}
		\end{tikzpicture}
		\end{center}
	\end{multicols}
	
		\vspace{5pt}
		\hrule
		\vspace{5pt}
		
	\begin{multicols}{2}
	Le translaté de $A+v$ par $v$ est $A+v+v = A+2v$, également représenté ci-contre.
	De façon similaire, on représente les translatés $A-v$, $A+1,5v$, $A - \frac{\pi}{6} v$ ci-dessous.
		\begin{center}
		\begin{tikzpicture}[>=stealth, scale=1]
		\begin{axis}[xmin = -10, xmax=10, ymin=-10, ymax=10, axis x line=none, axis y line=none, axis line style=<->, xlabel={}, ylabel={}, ticks = none]
			\addplot[BLUE_E, mark=*, mark size = 1] (4,2) node[above] {$A$};	
			% v = (2,-1)
			\addplot[very thick, ->, PINK] (0,0) -- (axis cs:8,-4) node[midway, below] {$v$};
			% A+v
			\addplot[dashed, ->, GREY] (4,2) -- (axis cs:12,-2) node[midway, below] {$v$};
			\addplot[RED_E, mark=*, mark size = 1] (12,-2) node[above right] {$A+v$};	
			% A + 2v
			\addplot[dashed, ->, GREY] (12,-2) -- (axis cs:20,-6) node[midway, below] {$v$};
			\addplot[GREEN_E, mark=*, mark size = 1] (20,-6) node[right] {$A+2v$};	
			% A - v
		\end{axis}
		\end{tikzpicture}
		\end{center}
	\end{multicols}
	
	
		\begin{center}
		\begin{tikzpicture}[>=stealth, scale=1]
		\begin{axis}[xmin = -10, xmax=10, ymin=-10, ymax=10, axis x line=none, axis y line=none, axis line style=<->, xlabel={}, ylabel={}, ticks = none]
			\addplot[BLUE_E, mark=*, mark size = 1] (4,2) node[above] {$A$};	
			% v = (2,-1)
			\addplot[very thick, ->, PINK] (0,0) -- (axis cs:8,-4) node[midway, below] {$v$};
			% A+v
			\addplot[dashed, ->, GREY] (4,2) -- (axis cs:12,-2) node[midway, below] {$v$};
			\addplot[RED_E, mark=*, mark size = 1] (12,-2) node[above right] {$A+v$};	
			% A + 2v
			\addplot[dashed, ->, GREY] (12,-2) -- (axis cs:20,-6) node[midway, below] {$v$};
			\addplot[GREEN_E, mark=*, mark size = 1] (20,-6) node[right] {$A+2v$};	
			% A - v
			\addplot[dashed, ->, GREY] (4,2) -- (axis cs:-4,6) node[midway, below] {$-v$};
			\addplot[GOLD_E, mark=*, mark size = 1] (-4,6) node[above right] {$A-v$};	
			% A + 1.5v
			\addplot[MAROON, mark=*, mark size = 1] (16,-4) node[above right] {$A+1,5v$};	
			% A - pi/6v
			\addplot[ORANGE, mark=*, mark size = 1] (4-3.14/6*8,2+3.14/6*4) node[above right] {$A-\frac{\pi}6v$};	
		\end{axis}
		\end{tikzpicture}
		\end{center}
	
	Remarquons que les points sont tous alignés.
	En ajoutant une infinité de tels points ; c'est-à-dire \emph{tous} les translatés de $A$ par un multiple de $v$, on obtient une droite.
	Cette droite ainsi créée est appelé la droite passant par $A$ et \emph{dirigée} par $v$, le vecteur $v$ donnant sa direction.
	
	En d'autres termes cette droite est l'ensemble
		\[ \bigset{ A+ k\cdot v \tq \text{ $k$ parcourt $\R$} }. \]
	
\section{Courbe représentative}
	
	Ça, vous savez !

\section{Montrer qu'un point appartient à une droite}

	Nommons $(d)$ la droite passant par $A$ et dirigée par $v$.
	Alors, un point $B$ appartient à cette droite si et seulement si $B$ est le translaté de $A$ par un multiple de $v$ :
		\[ B = A + k v, \]
	pour un $k\in\R$.
	
	Comme $k v$ envoie $A$ sur $B$, il est, par définition, le vecteur $\vec{AB}$.
	Dès lors, $v$ et $\vec{AB}$ sont colinéaires.
	
	Par conséquent, nous obtenons le
	\begin{lemma}
		$B$ appartient à la droite passant par $A$ et dirigée par $v$ si et seulement si $\vec{AB}$ et $v$ sont colinéaires.
	\end{lemma}

\end{document}
